\title{The Era of 1-bit LLMs: All Large Language Models are in 1.58 Bits}

\begin{abstract}
Emerging advancements, exemplified by BitNet~\cite{bitnet}, are ushering in the transformative epoch of 1-bit Large Language Models (LLMs). We present a novel 1-bit LLM variant, named \textbf{\text{BitNet b1.58}}, wherein each parameter (weight) adopts a ternary value set \{-1, 0, 1\}. This model achieves parity with full-precision Transformer LLMs (using FP16 or BF16) of equivalent architecture and training data volume, in terms of both perplexity and downstream task accuracy, while dramatically lowering costs relating to latency, memory footprint, computational throughput, and power usage. More importantly, the 1.58-bit LLM establishes a novel scaling law and methodology for constructing future LLMs that excel in both efficiency and performance. In addition, this innovation facilitates the emergence of a new computational paradigm and paves the way for custom hardware tailored for 1-bit LLMs.
\end{abstract}

\section{The Era of 1-bit LLMs}

Artificial intelligence has recently experienced rapid advances in the development and scale of Large Language Models (LLMs). These systems have reached outstanding benchmarks across a variety of natural language processing applications; however, their ever-growing parameter count presents notable difficulties for real-world deployment and accentuates environmental and financial challenges due to substantial energy demands.

A promising solution involves leveraging post-training quantization strategies to generate low-bit inference models~\cite{smoothquant, gptq, quip, quip_sharp}. By lowering the numerical precision of both weights and activations, this approach significantly decreases the memory load and computational cost associated with LLM operation. The prevailing movement has focused on pushing precision boundaries downward—from 16-bit representations to even 4-bit alternatives~\cite{gptq, awq}. Nevertheless, although post-training quantization is widely implemented throughout industry-grade LLMs, it remains a fundamentally sub-optimal compromise.

Recent advancements in 1-bit model designs, notably BitNet~\cite{bitnet}, suggest a compelling avenue for lowering the resource demands of LLMs while preserving their efficacy. Standard LLM implementations typically rely on 16-bit floating point formats (such as FP16 or BF16), and matrix multiplication dominates their computational workload. Consequently, the primary sources of computational expense are floating-point arithmetic operations. By contrast, BitNet’s matrix multiplications are limited to integer addition, which offers significant energy savings for LLMs. Since power consumption is a key constraint on performance in many hardware systems, these energy reductions can facilitate greater computational speed.

Beyond computation, moving model parameters from DRAM into the on-chip accelerator’s memory (for instance, SRAM) is a considerable contributor to inference overhead. While increasing SRAM capacity has been explored as a means to raise throughput, this approach incurs much higher costs compared to DRAM. 1-bit LLMs, however, require substantially less memory capacity and bandwidth than their full-precision counterparts, which can greatly decrease both the expenses and latency associated with loading weights from DRAM, leading to improved inference speed and efficiency.

This paper introduces a major new variant in the 1-bit LLM family, termed \textbf{\text{BitNet b1.58}}, where each model parameter is ternary, capable of taking the values \{-1, 0, 1\}. By introducing the additional value 0 to the classic 1-bit BitNet design, BitNet b1.58 utilizes 1.58 bits per parameter in binary encoding. \text{BitNet b1.58} maintains all the computational and memory benefits of the original BitNet, preserving its multiplication-free matrix multiplication approach that can be efficiently implemented. Its energy use remains comparable to the original BitNet, and it demonstrates markedly higher memory, throughput, and latency efficiency relative to FP16 LLM benchmarks. In addition, \text{BitNet b1.58} brings two key enhancements: first, it demonstrates improved modeling power by supporting explicit feature filtering via zero weights, enabling significant performance gains for 1-bit LLMs; and second, our empirical results indicate that, from the 3B parameter scale onward, \text{BitNet b1.58} achieves parity with full-precision (FP16) baselines in terms of both perplexity and downstream task outcomes, assuming equivalent settings for model size, training tokens, and related factors.

\section{BitNet b1.58}

\text{BitNet b1.58} is built upon the BitNet framework, which employs a Transformer architecture where the conventional \emph{nn.Linear} layers are substituted with \emph{BitLinear} layers. This model is trained from the ground up, utilizing 1.58-bit precision for its weights and 8-bit precision for its activations. In addition to the features of the original BitNet, several enhancements are incorporated, as described below.

\paragraph{Quantization Function.} To restrict the possible weight values to -1, 0, and +1, we use an \emph{absmean} quantization function. This approach initially normalizes the weight matrix by dividing it by the mean of its absolute values, and subsequently rounds each element to the closest integer in the set \{-1, 0, +1\}:

\begin{equation}
    \widetilde{W} = \mathrm{RoundClip}(\frac{W}{\gamma + \epsilon}, -1, 1),
\end{equation}

\begin{equation}
    \mathrm{RoundClip}(x, a, b) = \max(a, \min(b, \mathrm{round}(x))),
\end{equation}

\begin{equation}
    \gamma = \frac{1}{nm}\sum_{ij} |W_{ij}|.
\end{equation}

For the activation quantization, the method is consistent with what is used in BitNet. However, unlike the original, we do not normalize activations before non-linearities to fall within the interval $[0, Q_b]$. Instead, we linearly scale all activations to $[-Q_b, Q_b]$ on a per-token basis, eliminating the need for zero-point quantization. This strategy simplifies both software implementation and system-level optimization, while our experiments show only minimal impact on overall model performance.

\paragraph{LLaMA-alike Components.}

The LLaMA~\cite{llama, llama2} architecture has emerged as the foundational design for most open-source large language models. To align with community standards, our \text{BitNet b1.58} model adopts various LLaMA-alike architectural elements. In particular, it incorporates RMSNorm~\cite{rmsnorm}, SwiGLU~\cite{swiglu}, rotary embedding~\cite{rope}, and excludes all bias terms. This enables straightforward compatibility with leading open-source tools, such as Huggingface, vLLM~\cite{vllm}, and llama.cpp\footnote{\url{https://github.com/ggerganov/llama.cpp}}, requiring minimal adaptation effort.

\section{Results}

We benchmarked \text{BitNet b1.58} against our replicated FP16 \text{LLaMA LLM} models across multiple scales. To maintain fairness, both models underwent pre-training for 100 billion tokens using the RedPajama dataset~\cite{redpajama}. Zero-shot capabilities were assessed using several language understanding benchmarks, such as ARC-Easy~\cite{arc}, ARC-Challenge~\cite{arc}, Hellaswag~\cite{hellaswag}, Winogrande~\cite{winoGrande}, PIQA~\cite{piqa}, OpenbookQA~\cite{openbookqa}, and BoolQ~\cite{boolq}. We also included validation perplexity measurements on WikiText2~\cite{wikitext2} and C4~\cite{c4} to assess model generalization.

We further analyzed both \text{LLaMA LLM} and \text{BitNet b1.58} in terms of GPU memory consumption and inference latency. Experiments were conducted using the highly optimized FasterTransformer\footnote{\url{https://github.com/NVIDIA/FasterTransformer}} library for large language model inference on GPU. For \text{BitNet b1.58}, integration with the 2-bit kernel from Ladder~\cite{ladder} was used. Our primary latency metric was the duration required to generate each output token, reflecting inference efficiency.

In Table~\ref{tab:ppl}, we present a comparison of both perplexity and resource demands for \text{BitNet b1.58} versus \text{LLaMA LLM}. The findings indicate that at the 3B scale, \text{BitNet b1.58} achieves perplexity on par with the full precision \text{LLaMA LLM}, while offering a 2.71-fold improvement in speed and using only 28\% of the GPU memory (a 3.55$\times$ reduction). Notably, at 3.9B parameters, \text{BitNet b1.58} delivers 2.4$\times$ higher throughput and 3.32$\times$ less memory usage, outperforming the 3B \text{LLaMA LLM} variant by a significant margin.

Table~\ref{tab:zero-shot} details zero-shot accuracies on downstream tasks. For evaluation, we employed the standard \emph{lm-evaluation-harness}\footnote{\url{https://github.com/EleutherAI/lm-evaluation-harness}} protocol. The data show that with increasing scale, \text{BitNet b1.58} and \text{LLaMA LLM} demonstrate a converging accuracy gap. Crucially, from the 3B size upward, \text{BitNet b1.58} matches or exceeds the performance of its FP16 counterpart, consistent with our perplexity observations. In end-task accuracy, \text{BitNet b1.58} at 3.9B consistently surpasses \text{LLaMA LLM} 3B, requiring less memory and affording lower latency, illustrating \text{BitNet b1.58}'s Pareto dominance relative to state-of-the-art large language models.

\paragraph{Memory and Latency}

The model was also expanded to sizes of 7B, 13B, and 70B parameters, and corresponding computational costs were assessed. As depicted in Figure~\ref{fig:latency-memory-energy}, both latency and memory requirements display consistent scaling behaviors, with performance improvements becoming more pronounced as model size increases. Notably, at the 70B scale, \text{BitNet b1.58} demonstrates a speed-up of 4.1 times relative to the \text{LLaMA LLM} baseline. This acceleration arises from the scaling of \emph{nn.Linear} operation times, which increase with larger model architectures. Memory utilization follows an analogous trajectory: since embeddings are maintained in full precision, their proportionate contribution to overall memory footprint diminishes with larger models. All measurements were conducted using a 2-bit kernel, indicating that further efficiency gains are feasible with additional optimizations.

\paragraph{Energy}

Energy consumption attributable to arithmetic operations was estimated for both \text{BitNet b1.58} and \text{LLaMA LLM}, with a primary emphasis on matrix multiplication—the dominant contributor to computational costs in large language models. Figure~\ref{fig:energy} presents the energy breakdown, showing that \text{BitNet b1.58} relies mainly on INT8 additions, whereas \text{LLaMA LLM} employs a mix of FP16 additions and FP16 multiplications. Leveraging the energy estimation framework from~\cite{energycost, pokebnn}, we observe that \text{BitNet b1.58} achieves a 71.4-fold reduction in the arithmetic operations energy needed for matrix multiplication on 7nm process nodes. We also evaluated total energy usage in end-to-end runs processing 512 tokens, observing that efficiency gains in \text{BitNet b1.58} relative to the FP16 \text{LLaMA LLM} baseline become more significant as model size increases. This is primarily due to the rising proportion of the \emph{nn.Linear} component as models scale, whereas other cost sources contribute less at larger sizes.

\paragraph{Throughput}

We assess the throughput performance of both \text{BitNet b1.58} and the 70B-parameter \text{LLaMA LLM} by distributing them across two 80GB A100 GPUs and employing pipeline parallelism~\cite{gpipe} to facilitate running the \text{LLaMA LLM} 70B model on this hardware configuration. By progressively increasing the batch size until the available GPU memory was fully utilized (with a fixed sequence length of 512), we observe from Table~\ref{tab:throughput} that \text{BitNet b1.58} 70B accommodates a batch size up to 11 times greater than that of the \text{LLaMA LLM}, yielding an 8.9-fold boost in throughput.

\textbf{\text{BitNet b1.58} paves the way for a revised scaling law regarding the balance between model performance and inference efficiency}. Drawing upon Figures~\ref{fig:latency-memory-energy} and \ref{fig:energy}, we highlight the following correspondences between 1.58-bit and traditional 16-bit model sizes:

\begin{itemize}
    \item BitNet b1.58 at 13B parameters surpasses the 3B FP16 LLM in latency, memory demand, and energy efficiency.
    \item BitNet b1.58 at 30B parameters is superior to the 7B FP16 LLM in terms of latency, memory footprint, and energy consumption.
    \item BitNet b1.58 at 70B parameters offers improved latency, reduced memory usage, and lower energy consumption compared to the 13B FP16 LLM.
\end{itemize}

\paragraph{Training with 2T Tokens}

The volume of training tokens significantly impacts LLM performance. To evaluate how \text{BitNet b1.58} handles large-scale token counts, we trained a \text{BitNet b1.58} model using 2T tokens, adopting the same data curation strategy as StableLM-3B~\cite{StableLM-3B-4E1T}, which represents the leading open-source 3B model. Both our model and StableLM-3B were assessed across a suite of benchmarks—namely Winogrande~\cite{winoGrande}, PIQA~\cite{piqa}, SciQ~\cite{sciq}, LAMBADA~\cite{lambada}, and ARC-easy~\cite{arc}. We provide zero-shot accuracy results in Table~\ref{tab:2t}; for metrics based on accuracy and normalized accuracy, we report the mean value. The performance data for StableLM 3b at 2T tokens are referenced from its original report. Our results indicate that \text{BitNet b1.58} consistently outperforms across all benchmarked tasks, suggesting that LLMs with 1.58-bit quantization maintain strong generalization abilities.

\section{Discussion and Future Work}

\textbf{1-bit Mixture-of-Experts (MoE) LLMs}

The Mixture-of-Experts (MoE) paradigm has emerged as an efficient strategy for scaling LLMs while maintaining cost-effectiveness. Although this approach substantially lowers computation in terms of FLOPs, it is often hampered by elevated memory requirements and substantial inter-chip communication costs, limiting both its scalability and practical deployment. 1.58-bit LLMs offer promising solutions to these bottlenecks. The compact memory requirements of such models decrease the number of hardware devices necessary for running MoE architectures, thereby minimizing distributed system complexity. Furthermore, this compression notably lessens the cost associated with moving activations between devices. If a full model can be accommodated on a single chip, inter-chip communication overhead could be eliminated entirely.

\textbf{Native Support of Long Sequence in LLMs}

The capacity to process extended sequences is increasingly essential as LLM applications grow more sophisticated. A primary obstacle in long-sequence inference arises from the considerable memory used for KV caches. \text{BitNet b1.58} makes substantial progress in addressing this issue by quantizing activations from 16 bits down to 8 bits, thus effectively doubling the maximum context length for a fixed memory budget. This compression can be further optimized, with the possibility of achieving lossless storage at 4 bits or even less for 1.58-bit LLM variants, presenting a compelling direction for future investigation.

\textbf{LLMs on Edge and Mobile}

Deploying 1.58-bit LLMs stands to transform the utility of language models on edge and mobile hardware, where both memory and computational throughput are often severely constrained. By lowering energy usage and memory needs, 1.58-bit LLMs facilitate efficient inference in settings that previously could not support such models. This, in turn, expands the range of possible applications and enhances the capability of resource-limited platforms. Additionally, since these models are better suited to CPU architectures—which dominate edge and mobile ecosystems—\text{BitNet b1.58} can be run with greater efficiency and effectiveness in these contexts.

\textbf{New Hardware for 1-bit LLMs}

Recent innovations, such as Groq\footnote{\url{https://groq.com/}}, highlight the prospects for custom hardware, including Language Processing Units (LPUs), tailored for LLM workloads. Looking ahead, there is an evident opportunity to push this frontier by developing hardware and system architectures explicitly optimized for 1-bit LLM operations, motivated by the novel computation frameworks introduced by BitNet~\cite{bitnet}.